% Generated draft: requires mathematical and source review.
\documentclass[11pt]{article}
\usepackage[margin=25mm]{geometry}
\usepackage{fontspec}
\setmainfont{Latin Modern Roman}
\usepackage{amsmath,amssymb,mathtools,unicode-math}
\setmathfont{Latin Modern Math}
\usepackage{enumitem,needspace,xurl,hyperref}
\setlength{\parindent}{0pt}
\setlength{\parskip}{5pt}
\setlength{\emergencystretch}{3em}
\newcommand{\N}{\mathbb{N}}
\newcommand{\Z}{\mathbb{Z}}
\newcommand{\Q}{\mathbb{Q}}
\newcommand{\R}{\mathbb{R}}
\newcommand{\C}{\mathbb{C}}
\newcommand{\F}{\mathbb{F}}
\newcommand{\A}{\mathbb{A}}
\newcommand{\PP}{\mathbb{P}}
\newcommand{\E}{\mathbb{E}}
\newcommand{\eps}{\varepsilon}
\newcommand{\dd}{\,\mathrm{d}}
\begin{document}
\section*{UnsolvedMath ID: 700138}
\textbf{Source problem code:} catalogue entry
\subsection{Definitions and mathematical statement}
For a group \(G\), let \(\mathcal L(G)\) denote its subgroup lattice, ordered by inclusion. A lattice isomorphism \(\Phi:\mathcal L(G)\to\mathcal L(H)\) is a bijection such that, for all subgroups \(A,B\leq G\),
\[
A\leq B\quad\Longleftrightarrow\quad\Phi(A)\leq\Phi(B).
\]
Define the derived series by \(G^{(0)}=G\) and \(G^{(n+1)}=[G^{(n)},G^{(n)}]\). The group \(G\) is soluble if \(G^{(d)}=\{1\}\) for some nonnegative integer \(d\).

\textbf{Source question, with unresolved terminology.} For a pure soluble group \(G\), describe the lattice isomorphisms \(\mathcal L(G)\to\mathcal L(H)\), where \(H\) ranges over groups. Is every such \(G\) strictly determined by its lattice?

Here, interpreting ``lattice'' as the subgroup lattice is contextual. The supplied record defines neither ``pure soluble'' nor ``strictly determined''. Consequently, the class of admissible groups and the precise reconstruction requirement remain unspecified. In particular, the ordinary reconstruction assertion
\[
\forall H\;\bigl(\mathcal L(H)\cong\mathcal L(G)\Longrightarrow H\cong G\bigr),
\]
where \(H\) ranges over all groups, is a possible interpretation of determination by the subgroup lattice, but cannot be identified with the source's \emph{strict} determination without further evidence. It is also unresolved whether the question requires each individual lattice isomorphism to be induced by a group isomorphism.

The source states that the answer is affirmative for free soluble groups. This qualification is retained as a source assertion; the supplied record does not specify the meaning of strict determination or any rank or derived-length restrictions for that assertion.
\subsection{Short English statement}
The problem asks how subgroup-lattice isomorphisms behave for pure soluble groups and whether their lattices strictly determine the groups. The meanings of ``pure'' and ``strictly determined'' require clarification. The source reports an affirmative answer for free soluble groups.
\subsection{Sources}
\begin{itemize}
\item Supplied local UnsolvedMath registry record, permanent ID \(700138\), exact source problem code \texttt{catalogue entry}, titled \emph{Kourovka Notebook Problem 2.68}. Upstream revision: \texttt{d371ee917ac9647fbb06cd6fc529deec6f3dda82}.
\item The record attributes the problem to \emph{The Kourovka Notebook (Unsolved Problems in Group Theory)}, original problem 2.68, second issue (1967), as reproduced in the twenty-first issue edition (2026), arXiv:1401.0300v48. This bibliographic attribution is supplied evidence and has not been independently checked.
\item The proposer field is supplied as \texttt{L. E. Sadovski\textbackslash{}u\{\textbackslash{}i}; its final name encoding is incomplete. The supplied proposal year is 1967.
\item The collection description names E. I. Khukhro and V. D. Mazurov as editors of the notebook collection covering Issues 1--21 (1965--2026). The registry labels this record \emph{open}; this is a supplied status, not an independently verified mathematical or literature assessment.
\end{itemize}
\end{document}
